\section{Coarse-Graining and Emergence}

A self emerges. It is not a base thing. The base is the eight atoms, the mesh, the dock, the site, the tone, the lean, the knit, the wake, and the beat, and nothing more. So the live question is exactly where the self-level can live. The base knit makes bodies. The effective layer makes selves. The split between the two is not a matter of taste. It is a theorem.

\begin{principle}[The location of the self]
A reversible permutation has no degeneracy to coarse-grain, so a causal-emergent self-level cannot live in the bare knit. It lives at the layer where coarse-graining discards information.
\end{principle}

This section gives that result, the one requirement it rests on, and the transition matrix that measures it. Two terms recur, so a plain statement of each helps before we begin. The \emph{knit} is the law of how the state changes, collide then stream, a fixed reversible charge-conserving table on the $24$ sites of each dock followed by a one-step flow of every tone along its site. The knit is the base law, and it is the word we use for the state-law throughout, never any other. A \emph{self} is an emergent thing, a bound, self-correcting, identity-bearing knot woven across many docks over many beats. The base supplies the bodies. The effective layer supplies the selves.

\subsection{The one requirement}

Causal emergence needs \textbf{degeneracy}. Many micro causes funnel to one effect, so grouping the causes sharpens the macro into a stronger cause. For that to happen, information must be lost somewhere. There are only three places the loss can come from, and naming them fixes the whole problem.

\begin{table}[ht]
\centering
\begin{tabular}{@{}lll@{}}
\toprule
source & what it is & allowed in the base? \\
\midrule
the coarse map & the observer discards detail, a lossy observable on a reversible state & yes, the orthodox route, no base change \\
the dynamics & the knit itself discards detail, an attractor & no, this needs irreversibility the base forbids \\
the boundary & the system is open, the wake and the arrow & yes, an open growing mesh is not a closed permutation \\
\bottomrule
\end{tabular}
\caption{The three sources of the information loss that causal emergence requires.}
\label{tab:loss-sources}
\end{table}

Every route to a self-level is a way of supplying that loss from one of these three places. The coarse map is allowed and is the orthodox route. The boundary is allowed and is the wake. The dynamics route is the one barred from the base, because the base knit is reversible by construction, and that prohibition is the content of the next subsection.

\subsection{Reversibility forbids base-level causal emergence}

This is the core result. Causal emergence in the sense of Hoel is effective information rising under coarse-graining. It requires degeneracy in the micro dynamics, many causes funneling to one effect.

The bare knit is a \textbf{reversible permutation} on its full state. Every configuration of the mesh has exactly one successor under the knit and exactly one predecessor. A permutation has zero degeneracy and maximal effective information. So no coarse-graining of the true micro state can increase it.

\begin{theorem}[Reversibility blocks base-level emergence]
The base knit is a reversible permutation on the mesh state. A permutation has no degeneracy, so its effective information is already maximal and cannot be raised by coarse-graining the true micro state. Causal emergence is therefore impossible at the base. This is a theorem about reversibility, not an engineering limit.
\end{theorem}

It was measured directly, not merely argued. Over a long asymmetric orbit the micro successor map showed \textbf{zero injectivity violations} across $132$ distinct states, confirming that the knit is a permutation. Coarse-graining to the dock-charge field, the signed sum of the tones in a dock, compresses the state space by a factor of about $2.2$ and makes $48$ of $60$ coarse states non-deterministic. The information the coarse map throws away is exactly the prerequisite for emergence that the reversible base lacks.

\begin{result}[Loss is the prerequisite]
The micro successor map is injective, zero violations across $132$ states, so the base carries no degeneracy. The dock-charge coarse map compresses by $2.2$ and renders $48$ of $60$ coarse states non-deterministic. The degeneracy that emergence needs is created by the coarse map, never by the reversible base.
\end{result}

\subsection{The transition matrix}

The measurement runs on a transition matrix, the coarse map's dynamics written as probabilities. The construction is direct.

\begin{itemize}
\item Take a \textbf{lossy observable}, for example a block-averaged charge field. It is a many-to-one function of the true reversible state.
\item Run the reversible base knit beat by beat. Watch how the observable moves from one coarse state to the next.
\item Record the counts as a \textbf{Markov transition matrix}, the chance the macro lands in state $j$ given it was in state $i$.
\end{itemize}

A reversible micro turns into a probabilistic macro precisely because the coarse map threw detail away. The same micro state can sit under two different macro futures once you forget which micro state it was. Emergence appears when the macro has \textbf{metastable basins}, when many fine configurations funnel to one macro state. Then the macro is a sharper cause than any of its micro grains.

The engine measures real structure at the effective layer with this matrix, and each measure carries a control so the structure cannot be an artifact.

\begin{table}[ht]
\centering
\begin{tabular}{@{}lll@{}}
\toprule
measure & result & control \\
\midrule
second Markov eigenvalue, a slow mode & about $0.95$ & $0.15$ time-shuffled \\
slow mode across mesh sizes $48$, $64$, $96$ & survives & size sweep \\
implied timescale & rises then plateaus near $125$ beats & the standard validity check \\
learned effective knit commutes with one beat & far better than random & random-knit control \\
structured coarse map keeps effective information & more than random & random-map control \\
\bottomrule
\end{tabular}
\caption{Method checks at the effective layer, each with a control.}
\label{tab:effective-measures}
\end{table}

These are honest checks, not curve-fits. A slow mode at $0.95$ against a shuffled control of $0.15$ is a real timescale gap. An implied timescale that plateaus is the textbook validity test for a Markov-state model. A learned effective knit that commutes with one micro beat far better than a random one means the coarse dynamics is genuinely a coarse version of the same law. The effective level is well defined and measurable.

\subsection{The reversible bulk, the irreversible macro}

The architecture is the same one real physics uses, and stating that plainly removes the worry that something has gone wrong.

\begin{principle}[Reversible underneath, irreversible on top]
Reversible mechanics sits underneath. The second law and thermodynamics sit on top, recovered by coarse-graining. A reversible micro with an irreversible coarse macro is correct, not a bug.
\end{principle}

The realistic path to a clean emergent self-level is therefore not to make the base dissipative. It is to choose a better reversible knit, run it on the committed substrate, and do the emergence measurement between coarse levels. The options sort into faithful and forbidden.

\begin{table}[ht]
\centering
\begin{tabular}{@{}llll@{}}
\toprule
option & what changes & base-faithful & promise \\
\midrule
coarse-level emergence & analysis only & yes & needs a basin-forming dynamics \\
momentum-conserving collision & the knit, still reversible & yes & high, mobile interacting structures \\
richer sites, the $24$-direction $D_4$ & the substrate & yes & high, especially with the above \\
the wake and the arrow & the open mesh, the arrow emerges & yes & a real arrow, harder to measure \\
dissipative base & the base becomes irreversible & no & forbidden, a sixth ingredient \\
\bottomrule
\end{tabular}
\caption{Routes to an emergent self-level, faithful and forbidden.}
\label{tab:routes}
\end{table}

The last row is the line that must not be crossed. Dissipation is allowed and expected only at the declared effective layer. It is never put in the base. The wake is the legitimate boundary route, an open growing edge that supplies a genuine arrow without making the bulk irreversible.

\subsection{What the measurements found}

The momentum-conserving collision on the $24$ sites was run end to end. The mobility and scattering prerequisites pass. The clean emergent self-level is an honest negative, and that negative sharpens the boundary rather than blurring it.

\begin{table}[ht]
\centering
\begin{tabular}{@{}ll@{}}
\toprule
experiment & result \\
\midrule
mobile knit & pass, a lone particle travels ballistically, charge and momentum conserved \\
scatter & pass, two head-on particles are deflected onto a rotated axis, reversibly \\
coarse causal emergence, mobile gas & honest negative, no structured-over-random gain, no metastable mode \\
the wake and the arrow & pass, an open growing mesh carries a genuine arrow, echo $0.40$ open versus $0$ closed \\
soliton search & partial, free gliders are solitons, but two head-on gliders part rather than bind \\
\bottomrule
\end{tabular}
\caption{The end-to-end run of the momentum-conserving collision.}
\label{tab:findings}
\end{table}

The lesson is precise. Mobility and scattering are \textbf{necessary but not sufficient}. A causal-emergence gain needs a \textbf{metastable coarse mode}. The mobile gas lacks one, so its coarse map shows no structured-over-random gain. A flat self-diffusion observable did show the gain, because diffusion has the slow spatial mode the gas itself lacks. The wake supplies a real arrow with nothing added to the base. The closed bulk recovers under time reversal, an echo of $0$, while the open system does not, an echo of $0.40$, so the irreversibility lives only at the wake while the bulk stays reversible. Why the bulk is reversible and what the bath is, in full, is the section on the arrow and the bath.

The one remaining gap is \textbf{binding}, turning colliding solitons into lasting composites. The next move runs the mobile soliton dynamics inside the open growing mesh and tests whether colliding gliders are retained where the closed conservative knit let them part. A pass there is the wake turning solitons into selves.

\subsection{The conclusion}

The split is principled, and it can be stated in two lines.

\begin{result}[Bodies from the base, selves from the layer]
The base knit makes bodies. Persistent, confined, causally isolated structures arise from the bare reversible knit with no posited ingredients. The effective layer makes selves. Agency and causal emergence require information to be discarded, which the reversible base never does, so they are produced by coarse-graining, the information-losing step.
\end{result}

A self is therefore not hidden inside the base knit waiting to be found. It is what emerges when the base knit is coarse-grained. Reversibility is the reason, since a permutation preserves all information and so has no causal emergence. Coarse-graining is the cure, since it discards the information whose loss makes a macro-level a sharper cause than its micro grains.
